\documentclass[11pt]{article}
\usepackage[a4paper,margin=2.5cm]{geometry}
\usepackage[english]{babel}
\usepackage{amsmath,amssymb,amsthm}
\usepackage{libertinus}

\usepackage{fontawesome5}
\usepackage{xcolor}
\usepackage[most]{tcolorbox}

\usepackage{hyperref}
\hypersetup{ colorlinks=true,linkcolor=blue,urlcolor=cyan }

%\RequirePackage{luatex85}
\ProvidesPackage{mymath}

\RequirePackage{fontawesome5}
\RequirePackage{xcolor}
\RequirePackage[most]{tcolorbox}

\setlength{\parindent}{0pt}
\setlength{\parskip}{0.6em}

\tcbset{ thmstyle/.style={
    enhanced, breakable, colback=#1!5, colframe=#1!65!black, boxrule=0.5pt,
    arc=2pt, coltitle=white, fonttitle=\bfseries, attach boxed title to top
    left={yshift=-\tcboxedtitleheight/2, xshift=3mm}, boxed title
    style={colback=#1!65!black, arc=1pt, boxrule=0pt}, top=4mm, bottom=2mm,
    before skip=1.4em, after skip=1em,
} }

\newtcbtheorem[number within=section]{definition}{Definition}%
{thmstyle=green}{def}
\newtcbtheorem[use counter from=definition]{lemma}{Lemma}%
{thmstyle=orange}{lem}
\newtcbtheorem[use counter from=definition]{proposition}{Proposition}%
{thmstyle=blue}{prop}
\newtcbtheorem[use counter from=definition]{theorem}{Theorem}%
{thmstyle=red}{thm}

\newtcolorbox{important}{
    enhanced, breakable, frame hidden, borderline west={2.5pt}{0pt}{violet!70},
    colback=violet!6, colbacktitle=violet!6, coltitle=violet!70!black,
    fonttitle=\bfseries, title={\faExclamationCircle\ Important}, before
    skip=1em, after skip=1em,
}
\newtcolorbox{warning}{
    enhanced, breakable, frame hidden, borderline
    west={2.5pt}{0pt}{orange!70!black}, colback=orange!8, colbacktitle=orange!8,
    coltitle=orange!60!black, fonttitle=\bfseries,
    title={\faExclamationTriangle\ Warning}, before skip=1em, after skip=1em,
}

\usepackage{bm}

\newenvironment{tightcenter}{%
    \setlength\topsep{0pt}
    \setlength\parskip{0pt}
    \begin{center}
}{%
    \end{center}
}

\title{Stupid penduings}
\author{BadDev8}

\begin{document}
\maketitle

\section*{Introduction}

This doc will brefly explain what the algorithm used in this project is and how
it works.

\section*{Basics}

\begin{warning}
    Notice that we assume all values follow a normal / gaussian distribution.

    This is not necessarily true, but it is "precise enough" for the classifier
    to get 97\% (or more) precision.
\end{warning}

The classifier itself is, in fact, really simple. \\
It's called \textbf{Naive Bayesian classifier} (or NBC) and it's based on the
bayes formula:

\begin{proposition}
    {Bayes formula}{} Given two possible events A and B, it holds:
    \[
        P(A|B) = \frac{P(B|A) P(A)}{P(B)}
    \]
    where $P(A|B)$ is the probability of the event A given that event B holds.
\end{proposition}

In the training, the classifier is given examples of instances, each one with a
label corresponding to his class. So, when it has to classify a new instance,
the classifier calculates the probability of the instance to be in every class,
and chooses the class with the heighest probability.

In this specific dataset our instances are penguins and our classes are species.
This means that the NBC wants to calculate the probabilities:

\[
    P(Specie|x_1,x_2, \dots \ , x_n) = \frac{P(x_1,x_2, \dots \ , x_n|Specie) P(Specie)}{P(x_1,x_2, \dots \ , x_n)}
\]

\begin{tightcenter}
    where $x_1,x_2, \dots , x_n$ are some data of a specific penguin (bill
    lenght, flipper lenght, \dots)
\end{tightcenter}
for each specie (Adelie, Chinstrap, Gentoo), then compare them and choose the
specie with the higher probability. \\
\\
Notice that $P(x_1,x_2, \dots \ , x_n)$ doesn't really need to be evaluated,
since dividing two number for the same positive value does not change their
order. So we can just ignore it and get the faster:
\[
    P(x_1,x_2, \dots \ , x_n|Adelie) P(Adelie)
\]

Let's look to the specific components of this formula:
\begin{itemize}
    \item $P(Adelie)$ is really easy to get, since it is the same as the
          frequency of the specie divided by the number of total samples.
    \item $P(x_1,x_2, \dots \ , x_n|Adelie)$ is tougher and needs a chapter for
          itself.
\end{itemize}

\section*{How do we calculate that?}

$P(x_1,x_2, \dots , x_n|Adelie)$ is a distribution on n variables. Notice that
there are ways to calculate this type of distribution and the classifiers who do
that are called \textbf{Bayes Optimal Classifier}. \\
They reach the best accuracy
(\href{https://en.wikipedia.org/wiki/Bayes_classifier}{literally}) but can be
incredibly expensive from the computational point of view.

This is why we often assume that the single variables are independent between
them. This is not necessarily true and can lead to errors, but at the same time
it's ridicoluosly good to make computation faster.

Assuming independencies, we can then rewrite the formula as:

\[
    P(x_1,x_2, \dots \ , x_n|Adelie) = P(x_1|Adelie) \cdot P(x_2|Adelie) \cdot \dots \cdot P(x_n|Adelie)
\]

which are probabilities on mono-variable normal distributions.

This can be calculated easily knowing the distributions itself. The process to
get this type of distributions is easy, but understanding it requires some
decent calculus knowledge, which can be explored in the Appendix.

I'm going to describe the process on only one of the variables, since we can
repeat the same process for the others.\\
We start by removing that "given Adelie" (or the other specie we are
considering) by just \dots using only the samples from Adelie penguins!

We need to use this important lemma:
\begin{lemma}{}{}
    A normal distribution on one varible is univocally defined by two values
    $\overline{x}, \sigma$ (which represent the mean and the standard deviation)
    as the function:
    \[
        f(x) = \frac{1}{\sigma \sqrt{2\pi}} \ \ e^{- \frac{(x-\overline{x})^2}{2\sigma^2}}
    \]
\end{lemma}

notice that those values are ideal, we can't have them. But we can have a good
estimation of them through the mean and the standard deviation of the data we
have.

Usually to slightly improve performance we use another fact about the average
error from the mean (MEA):
\begin{proposition}{}{}
    In a gaussain distribution with one varible it holds:
    \[
        \text{MAE} = \sqrt{\frac{2}{\pi}} \ \sigma
    \]
\end{proposition}
Calculations for the MAE is a bit faster than those fot the standard deviation,
so this can be used to improve performance.

\section*{Conclusion}

Repeating this process over and over we can have a distribution for every
element of every pair (specie, variable). We can then calculate the three
probalities using:

\[
    P(x_1,x_2, \dots \ , x_n|Adelie) P(Adelie)
\]

and chose the specie with the higher probability.

\section*{Appendix}

Here are presented demonstations for some facts used in the exposing.

\begin{proposition}{}{}
    A normal distribution on one varible is univocally defined by two values
    $\overline{x}, \sigma$ (which represent the mean and the standard deviation)
    as the function:
    \[
        f(x) = \frac{1}{\sigma \sqrt{2\pi}} \ \ e^{- \frac{(x-\overline{x})^2}{2\sigma^2}}
    \]
\end{proposition}

\begin{proof}
    The normal distribution is defined from the function:
    \[
        f(x) = \text{A} e^{-\text{h}^2(x-\text{m})^2}
    \]
    where
    \begin{itemize}
        \item A is the absolute maxsimum of the function
        \item m is the point where the function has its max value A
        \item h is a parameter for how "flat" the curve is
    \end{itemize}

    It has to hold:
    \[
        \int_{-\infty}^{\infty} f(x) dx = 1
    \]
    from which:
    \[
        \int_{-\infty}^{\infty} \text{A} e^{-\text{h}^2(x-\text{m})^2} dx = 1
    \]
    changing $h(x - m) = z$
    \[
        \text{A} \int_{-\infty}^{\infty} e^{-z^2} dz = \text{h}
    \]
    but since we know $\int_{-\infty}^{\infty} e^{-z^2} = \sqrt{\pi}$ it holds:
    \[
        \text{A} = \frac{\text{h}}{\sqrt{\pi}}
    \]

    this means we can rewrite everything into:
    \[
        f(x) = \frac{\text{h}}{\sqrt{\pi}} e^{-\text{h}^2(x-\text{m})^2}
    \]

    Now, the mean $\overline{x}$ is defined in continuos distributions as:
    \[
        \overline{x} = \int_{-\infty}^{\infty} x f(x) dx
    \]
    from which with the same substitution we get:
    \[
        \overline{x} = \frac{\text{h}}{\sqrt{\pi}} \int_{-\infty}^{\infty} x e^{-\text{h}^2(x-\text{m})^2} dx
    \]
    \[
        \overline{x} = \frac{}{\sqrt{\pi}} \int_{-\infty}^{\infty} (\frac{z}{\text{h}} + \text{m}) e^{-z^2} dz
    \]
    the first term of the integral is 0 since it's an odd function, it remains:
    \[
        \overline{x} = \frac{1}{\sqrt{\pi}} \int_{-\infty}^{\infty} \text{m} e^{-z^2} dz
                     = \frac{1}{\sqrt{\pi}} \ \text{m} \ \sqrt{\pi}
                     = m
    \]

    In a similar way we can get the Square error $\sigma^2$:
    \[
        \sigma^2 = \int_{-\infty}^{\infty} (x - \overline{x})^2 f(x) dx
    \]
    \[
        \sigma^2 = \int_{-\infty}^{\infty} \frac{\text{h}}{\sqrt{\pi}} (x - m)^2 e^{-\text{h}^2(x-\text{m})^2} dx
    \]
    \[
        \sigma^2 = \frac{1}{\sqrt{\pi} \ \text{h}^2} \int_{-\infty}^{\infty} (z)^2 e^{-z^2} dz
    \]
    using integration by part we get:
    \[
        \int_{-\infty}^{\infty} z^2 e^{-z^2} = \frac{\sqrt{\pi}}{2}
    \]
    and:
    \[
        \sigma^2
        =           \frac{1}{\sqrt{\pi} \ \text{h}^2} \ \frac{\sqrt{\pi}}{2}
        =           \frac{1}{2 \ \text{h}^2}
        \rightarrow \sigma
        =           \sqrt{\frac{1}{2 \ \text{h}^2}}
        =           \frac{1}{\sqrt{2} \ \text{h}}
    \]
    \[
        \text{h} = \frac{1}{\sqrt{2} \sigma}
    \]

    By substituting into the original equation we found the thesis:
    \[
        f(x) = \frac{\text{h}}{\sqrt{\pi}} e^{-\text{h}^2(x-\text{m})^2}
             = \frac{1}{\sqrt{2 \pi \sigma}} e^{-(x-\overline{x})^2 / 2\sigma^2}
    \]

\end{proof}

\begin{proposition}{}{}
    In a gaussain distribution with one varible it holds:
    \[
        \text{MAE} = \sqrt{\frac{2}{\pi}} \ \sigma
    \]
\end{proposition}
\begin{proof}
    By definition:
    \[
        \text{MAE}
        = \int_{-\infty}^{\infty} |x
          - \text{m}| \frac{\text{h}}{\sqrt{\pi}} e^{-\text{h}^2(x-\text{m})^2} dx
    \]
    with $\text{y} = |x - \text{m}|$
    \[
        \text{MAE} = \frac{2\text{h}}{\sqrt{\pi}} \int_{0}^{\infty} \text{y} e^{-\text{h}^2 \text{y}^2} dy
    \]
    with
    \[
        \text{z} = \text{h}^2 \text{y}^2
    \]
    \[
        dz = 2 \text{h}^2 \text{y} \ dy
    \]
    \[
        \text{y} \ dy = \frac{dz}{2 \text{h}^2}
    \]
    we get:
    \[
        \text{MAE} = \frac{1}{\text{h}\sqrt{\pi}} \int_{0}^{\infty} e^{-\text{z}} dz
    \]
    \[
        \text{MAE} = \frac{1}{\text{h}\sqrt{\pi}}
    \]

    But from the prevous proposition we know
    $\text{h} = \frac{1}{\sqrt{2} \sigma}$ so:
    \[
        \text{MAE} = \frac{1}{\text{h}\sqrt{\pi}}
                   = \frac{\sqrt{2} \sigma}{\sqrt{\pi}}
                   = \sqrt{\frac{2}{\pi}} \ \sigma
    \]

\end{proof}

\end{document}
